\begin{solapartat}
\punts{0,5} Segons la llei de la gravitació universal, el mòdul de la força sobre el satèl·lit per l'atracció de la Terra és:
\[ F = G\,\frac{m_s M_T}{r^2} \]
La segona llei de Newton estableix que: $\vec{F} = m_s\vec{a}$

D'altra banda, considerant que el satèl·lit descriu un moviment circular uniforme al voltant de la Terra, la seva acceleració és l'acceleració centrípeta: $a = \frac{v^2}{r} = \omega^2 r$.

Com que sobre el satèl·lit només actua la força de la gravetat:
\[ G\,\frac{m_s M_T}{r^2} = \frac{m_s v^2}{r} \;\Rightarrow\; v = \sqrt{G\,\frac{M_T}{r}} \]

\punts{0,25} El radi $r = R_T + h = (6370 + 550)\ \mathrm{km} = 6,92\times 10^{6}\ \mathrm{m}$.

Utilitzant l'expressió, obtenim el valor de la velocitat orbital per al satèl·lit:
\[ v = \sqrt{G\,\frac{M_T}{r}} = \sqrt{6,67\times 10^{-11}\,\frac{5,98\cdot 10^{24}}{6,92\cdot 10^{6}}} = 7,59\times 10^{3}\ \mathrm{m/s} \]

\punts{0,5} El nombre de voltes al voltant de la Terra que fa el satèl·lit durant un dia. Hem de comparar els dos temps.

El període del satèl·lit: $T = \dfrac{2\pi r}{v} = \dfrac{2\pi\, 6,92\cdot 10^{6}}{7,59\cdot 10^{3}} = 5,73\times 10^{3}\ \mathrm{s}$

I un dia amb segons: $t_{dia} = 24\ \mathrm{h}\,\dfrac{3600\ \mathrm{s}}{1\mathrm{h}} = 8,64\times 10^{4}\ \mathrm{s}$

I el nombre de voltes: $n_{\frac{voltes}{dia}} = \dfrac{t_{dia}}{T} = \dfrac{86,4\times 10^{3}}{5,73\times 10^{3}} = 15,1\ \text{voltes}$
\end{solapartat}

\begin{solapartat}
\punts{0,75} L'energia mecànica és la suma de l'energia cinètica i potencial:
\[ E_m = E_c + E_p = \frac{1}{2} m_s v^2 - G\,\frac{m_s M_T}{r} = \frac{1}{2} m_s G\,\frac{M_T}{r} - G\,\frac{m_s M_T}{r} = -\frac{1}{2}\,G\,\frac{m_s M_T}{r} \]

\punts{0,25} Calculem l'energia mecànica a l'òrbita de 530~km d'alçada:

radi $r = 6,37\times 10^{6} + 5,3\cdot 10^{5} = 6,90\times 10^{6}\ \mathrm{m}$

I utilitzant l'expressió, obtenim el valor de l'energia mecànica per al satèl·lit:
\[ E_m = -\frac{1}{2}\,G\,\frac{m_s M_T}{r} = -\frac{1}{2}\,6,67\times 10^{-11}\,\frac{227\times 5,98\times 10^{24}}{6,90\times 10^{6}} = -6,56\times 10^{9}\ \mathrm{J} \]

\punts{0,25} El signe negatiu de l'energia mecànica indica que el satèl·lit està seguint una òrbita tancada (o lligada) al voltant de la Terra.
\end{solapartat}
